\documentclass[11pt]{article}

\usepackage{amsmath}
\usepackage{amssymb}
\usepackage{booktabs}
\usepackage{enumitem}
\usepackage[margin=1in]{geometry}

\begin{document}

\begin{flushright}
CSE 21 Fall 2026\\
Homework 7\\
Submitted: 11/23/2026
\end{flushright}

\begin{center}
David Washington(200484943), Judah Boyce(200453537),\\
Osinachi Agaranna(200455712), Nick Young(200455733)
\end{center}

\section*{Problem 1}

\subsection*{(a)}

Any domino tiling of a $2\times 2n$ grid falls into one of two cases, based on the vertical line between column $n$ and column $n+1$.

\textbf{Case 1: no domino crosses the middle.}
The tiling splits into a tiling of the left $2\times n$ grid and a tiling of the right $2\times n$ grid. These choices are independent, so there are $DT(n)^{2}$ tilings in this case.

\textbf{Case 2: a domino crosses the middle.}
A horizontal domino covers two consecutive columns of one row. In a $2$-row tiling, horizontal dominos come in stacked pairs: if one row has a horizontal domino on columns $i$ and $i+1$, the other row must have one on the same two columns, because a vertical domino needs both rows of a single column. So a domino crossing the middle forces a block of two horizontal dominos covering columns $n$ and $n+1$. Removing that block leaves a $2\times(n-1)$ grid on the left and a $2\times(n-1)$ grid on the right, tiled independently. There are $DT(n-1)^{2}$ tilings in this case.

The two cases are disjoint, and every tiling of the $2\times 2n$ grid belongs to one of them: the middle line is either crossed by a domino or it is not. Therefore
\[
DT(2n)=DT(n)^{2}+DT(n-1)^{2}.
\]

\[
\boxed{DT(2n)=DT(n)^{2}+DT(n-1)^{2}}
\]

\subsection*{(b)}

Fix an integer $n\ge 1$, and consider a domino tiling of a $2\times(2n+1)$ grid. The width $2n+1$ is odd. Horizontal dominos cover two columns at a time, so they cover an even number of columns altogether. Hence every such tiling uses at least one vertical domino, and there is a leftmost vertical domino.

That leftmost vertical domino lies in an odd-numbered column. If it lay in column $2k$ for some $k\ge 1$, then columns $1$ through $2k-1$ would be tiled with no vertical domino. That region has odd width, and horizontal pairs cover two columns at a time, so it cannot be tiled. The possible columns for the leftmost vertical domino are therefore
\[
1,3,5,\ldots,2n+1.
\]

\textbf{The leftmost vertical domino is in column $2j+1$, for some $j$ with $0\le j\le n$.}
Columns $1$ through $2j$ contain no vertical domino, so they are filled by stacked horizontal pairs. A horizontal domino covers two consecutive columns of one row and must be stacked on the other row, so the only such tiling is a stacked pair on columns $1$ and $2$, then columns $3$ and $4$, and so on through columns $2j-1$ and $2j$. A $2\times 2j$ grid therefore has exactly one tiling with no vertical domino. Column $2j+1$ is the vertical domino itself. Removing those $2j+1$ columns leaves columns $2j+2$ through $2n+1$, a grid of width
\[
(2n+1)-(2j+1)=2(n-j).
\]
That remaining grid may be tiled in any way, so it contributes $DT\bigl(2(n-j)\bigr)$ tilings. The left side was forced, so the whole case has $DT\bigl(2(n-j)\bigr)$ tilings.

The two smallest cases make the pattern explicit.

\begin{itemize}[leftmargin=1.6em, itemsep=0.35em, topsep=0.4em]
\item If the leftmost vertical domino is in column $1$, then $j=0$. Removing it leaves a $2\times 2n$ grid, which has $DT(2n)$ tilings.
\item If the leftmost vertical domino is in column $3$, then $j=1$. Columns $1$ and $2$ are the unique stacked pair of horizontal dominos. Removing those two columns and the vertical domino leaves a $2\times(2n-2)$ grid, which has $DT(2n-2)$ tilings.
\end{itemize}

Running through $j=0,1,\ldots,n$ gives the terms
\[
DT(2n)+DT(2n-2)+\cdots+DT(2)+DT(0).
\]
When $j=n$, the vertical domino is in column $2n+1$ and nothing remains to its right, which is the empty tiling counted by $DT(0)=1$.

Every tiling has exactly one leftmost vertical domino, so these cases are disjoint and they cover every tiling of the $2\times(2n+1)$ grid. Therefore
\[
DT(2n+1)=DT(0)+DT(2)+DT(4)+\cdots+DT(2n).
\]

\[
\boxed{DT(2n+1)=DT(0)+DT(2)+\cdots+DT(2n)}
\]

\section*{Problem 2}

\subsection*{(a)}

\textbf{(i)}
A nonempty string over $\{1,4,5\}$ ends in $1$, $4$, or $5$. Deleting that last digit leaves a string whose digits sum to $n-1$, $n-4$, or $n-5$, and every such string arises once. Thus, for $n\ge 5$,
\[
A(n)=A(n-1)+A(n-4)+A(n-5).
\]
The empty string is the one string that sums to $0$, so $A(0)=1$. The strings for the remaining base cases are:
\begin{align*}
A(1)&=1 &&\text{(``$1$'')},\\
A(2)&=1 &&\text{(``$11$'')},\\
A(3)&=1 &&\text{(``$111$'')},\\
A(4)&=2 &&\text{(``$1111$'', ``$4$'')}.
\end{align*}
The same recurrence, with $A(m)=0$ for $m<0$, also produces these four values. As a check, $A(5)=A(4)+A(1)+A(0)=2+1+1=4$, and the strings are ``$11111$'', ``$14$'', ``$41$'', and ``$5$''.

\[
\boxed{A(n)=A(n-1)+A(n-4)+A(n-5)\ (n\ge 5),
\quad A(0)=A(1)=A(2)=A(3)=1,\ A(4)=2}
\]

\textbf{(ii)}
Substitute $A(n)=x^{n}$ into the recurrence and divide by $x^{n-5}$:
\[
x^{5}=x^{4}+x+1, \qquad x^{5}-x^{4}-x-1=0.
\]
Wolfram gives one real root $r\approx 1.49709$. The other four roots come in conjugate pairs, and their absolute values are about $1.098$ and $0.744$. Both are smaller than $r$, and $r$ is a simple root, so
\[
A(n)=c\,r^{n}+\text{(terms from the smaller roots)}.
\]
This is $\Theta(r^{n})$ once $c\neq 0$.

Strings that use only the digits $1$ and $4$ are part of the count $A(n)$. Let $C(n)$ be the number of those strings. The last digit is $1$ or $4$, so $C(n)=C(n-1)+C(n-4)$ with $C(0)=1$. The characteristic polynomial is $x^{4}-x^{3}-1=0$. Wolfram gives a real root $s\approx 1.380$, and the other three roots have absolute value less than $1$. Since $C(n)\ge 1$, the $s^{n}$ term is present and $C(n)=\Theta(s^{n})$.

Every such string is also counted by $A(n)$, so $A(n)\ge C(n)$. If $c$ were $0$, the largest roots left in $A(n)$ would have absolute value about $1.098$, and $A(n)=O(1.098^{n})$. That growth is smaller than $\Theta(s^{n})$, because $s\approx 1.380>1.098$. So $c\neq 0$, and
\[
A(n)=\Theta(r^{n}).
\]

\[
\boxed{A(n)=\Theta(r^{n}),\quad r\approx 1.49709}
\]

\subsection*{(b)}

\textbf{(i)}
Let $T_{i}(n)$ be the number of valid length-$n$ strings that start with $i$. Partitioning on the first character gives
\[
T(n)=T_{0}(n)+T_{1}(n)+T_{2}(n).
\]
A string starting with $2$ can be followed by any valid string of length $n-1$: the digit $2$ does not form $00$ or $11$ with whatever comes next. So $T_{2}(n)=T(n-1)$. Swapping $0$ and $1$ is a bijection between the valid strings starting with $0$ and those starting with $1$, so $T_{0}(n)=T_{1}(n)$. Therefore
\[
T(n)=2T_{0}(n)+T(n-1), \qquad T_{0}(n)=\frac{T(n)-T(n-1)}{2}.
\]
For a second expression, drop the leading $0$. The remaining string is any valid string of length $n-1$ that does not itself start with $0$, since a leading $00$ is forbidden, and every such string stays valid after a $0$ is written in front. So
\[
T_{0}(n)=T(n-1)-T_{0}(n-1).
\]
The same formula one step earlier is $T_{0}(n-1)=(T(n-1)-T(n-2))/2$. Equating the two expressions for $T_{0}(n)$ gives
\begin{align*}
\frac{T(n)-T(n-1)}{2}
&= T(n-1)-\frac{T(n-1)-T(n-2)}{2}.
\end{align*}
Multiplying through by $2$ and simplifying,
\begin{align*}
T(n)-T(n-1)
&= 2T(n-1)-\bigl(T(n-1)-T(n-2)\bigr)\\
&= T(n-1)+T(n-2),
\end{align*}
so $T(n)=2T(n-1)+T(n-2)$ for $n\ge 2$.

The empty string is the one string of length $0$, and it has no forbidden pair, so $T(0)=1$. The length-$1$ strings are $0$, $1$, and $2$, so $T(1)=3$.

\[
\boxed{T(n)=2T(n-1)+T(n-2)\ (n\ge 2),\quad T(0)=1,\quad T(1)=3}
\]

\textbf{(ii)}
The prompt asks for a bound on a sequence $B(n)$. That sequence is $T(n)$ from part (i). The characteristic equation of $T(n)-2T(n-1)-T(n-2)=0$ is
\[
x^{2}-2x-1=0, \qquad x=1\pm\sqrt{2}.
\]
Numerically, $1+\sqrt{2}\approx 2.414$ and $1-\sqrt{2}\approx -0.414$. The general solution is
\[
T(n)=c_{1}(1+\sqrt{2})^{n}+c_{2}(1-\sqrt{2})^{n}.
\]
The initial conditions are $c_{1}+c_{2}=1$ and $c_{1}(1+\sqrt{2})+c_{2}(1-\sqrt{2})=3$. Substitute $c_{2}=1-c_{1}$:
\begin{align*}
c_{1}(1+\sqrt{2})+(1-c_{1})(1-\sqrt{2})&=3,\\
c_{1}\bigl((1+\sqrt{2})-(1-\sqrt{2})\bigr)&=3-(1-\sqrt{2}),\\
2\sqrt{2}\,c_{1}&=2+\sqrt{2},\\
c_{1}&=\frac{2+\sqrt{2}}{2\sqrt{2}}=\frac{1+\sqrt{2}}{2}\neq 0.
\end{align*}
Since $|1-\sqrt{2}|<1<1+\sqrt{2}$, the second term is bounded and the first term determines the growth:
\[
T(n)=\Theta\!\left((1+\sqrt{2})^{n}\right).
\]

\[
\boxed{T(n)=\Theta\!\left((1+\sqrt{2})^{n}\right)}
\]

\section*{Problem 3}

\subsection*{(a)}

$P(0)$ is a plus with four arms, so it has four endpoints: $S(0)=4$. The prompt states the recursive step for $n\ge 2$. The picture builds $P(1)$ from $P(0)$ by that same step, so the relation below holds for every $n\ge 1$. At step $n$, a T is attached at every endpoint of $P(n-1)$. In the picture each T is a stem with a two-sided crossbar, and the endpoint counts are $S(0)=4$, $S(1)=12$, and $S(2)=36$. Those are $4$, $4\cdot 3$, and $4\cdot 3^{2}$, so each old endpoint is replaced by three new endpoints:
\[
S(n)=3S(n-1).
\]

\[
\boxed{S(n)=3S(n-1),\quad S(0)=4}
\]

\subsection*{(b)}

Substitute the recurrence from (a) over and over. Each time, $S(m)$ is replaced by $3S(m-1)$, so the power of $3$ goes up by one and the index goes down by one:
\[
S(n)=3S(n-1)=3^{2}S(n-2)=\cdots=3^{n}S(0).
\]
After $n$ steps the index is $0$. Using $S(0)=4$,
\[
S(n)=3^{n}\cdot 4=4\cdot 3^{n}.
\]

\[
\boxed{S(n)=4\cdot 3^{n}}
\]

\subsection*{(c)}

The original plus is four arms of length $1$, so $L(0)=4$. At step $n\ge 1$, a T is added at each of the $S(n-1)$ endpoints. Each T contributes three arms, and each of those arms has length equal to half the length of the arms added at the previous step, namely $(1/2)^{n}$. Segments added at earlier steps keep their original lengths. The number of new segments is
\[
3S(n-1)=3\cdot 4\cdot 3^{n-1}=4\cdot 3^{n}.
\]
The length added at step $n$ is
\[
4\cdot 3^{n}\cdot\left(\frac{1}{2}\right)^{n}=4\left(\frac{3}{2}\right)^{n}.
\]
Therefore
\[
L(n)=L(n-1)+4\left(\frac{3}{2}\right)^{n}.
\]

\[
\boxed{L(n)=L(n-1)+4\left(\dfrac{3}{2}\right)^{n},\quad L(0)=4}
\]

\subsection*{(d)}

Unfold the recurrence from (c) down to $L(0)$. The value $L(n)$ is the original length plus the length added at each later step:
\[
L(n)=L(0)+4\cdot\frac{3}{2}+4\left(\frac{3}{2}\right)^{2}+\cdots+4\left(\frac{3}{2}\right)^{n}.
\]
Since $L(0)=4=4\cdot(3/2)^{0}$, every term has the same shape:
\[
L(n)=\sum_{k=0}^{n}4\left(\frac{3}{2}\right)^{k}.
\]
The sum $\sum_{k=0}^{n}r^{k}=(r^{n+1}-1)/(r-1)$ with $r=3/2$ gives
\begin{align*}
L(n)
&=4\cdot\frac{\left(\dfrac{3}{2}\right)^{n+1}-1}{\dfrac{3}{2}-1}
=4\cdot\frac{\left(\dfrac{3}{2}\right)^{n+1}-1}{\dfrac{1}{2}}
=8\left(\left(\frac{3}{2}\right)^{n+1}-1\right).
\end{align*}
The first term simplifies by $8\cdot(3/2)^{n+1}=8\cdot(3/2)\cdot(3/2)^{n}=12\cdot(3/2)^{n}$, so
\[
L(n)=12\left(\frac{3}{2}\right)^{n}-8.
\]

\[
\boxed{L(n)=12\left(\dfrac{3}{2}\right)^{n}-8}
\]

\section*{Problem 4}

Let $\Sigma=\{A,B,C,D,E,F,G,H,I,J,K,L\}$, so $|\Sigma|=12$.

\subsection*{(a)}

A fixed-length code for one character needs the smallest $k$ with $2^{k}\ge 12$. Since $2^{3}=8<12$ and $2^{4}=16\ge 12$, each character uses $4$ bits. A string of $5$ characters therefore uses
\[
4\cdot 5=20
\]
bits.

\[
\boxed{20}
\]

\subsection*{(b)}

Number the letters in order, starting at $0$, and write each number as a $4$-bit string.

\begin{center}
\begin{tabular}{ccc}
\toprule
Letter & Decimal & Binary \\
\midrule
$A$ & $0$ & $0000$ \\
$B$ & $1$ & $0001$ \\
$C$ & $2$ & $0010$ \\
$D$ & $3$ & $0011$ \\
$E$ & $4$ & $0100$ \\
$F$ & $5$ & $0101$ \\
$G$ & $6$ & $0110$ \\
$H$ & $7$ & $0111$ \\
$I$ & $8$ & $1000$ \\
$J$ & $9$ & $1001$ \\
$K$ & $10$ & $1010$ \\
$L$ & $11$ & $1011$ \\
\bottomrule
\end{tabular}
\end{center}

\subsection*{(c)}

\textbf{(i)}
\begin{align*}
B&\mapsto 0001, &
E&\mapsto 0100, &
L&\mapsto 1011, &
C&\mapsto 0010, &
H&\mapsto 0111.
\end{align*}
\[
\boxed{0001\,0100\,1011\,0010\,0111}
\]

\textbf{(ii)}
\begin{align*}
L&\mapsto 1011, &
I&\mapsto 1000, &
K&\mapsto 1010, &
E&\mapsto 0100, &
D&\mapsto 0011.
\end{align*}
\[
\boxed{1011\,1000\,1010\,0100\,0011}
\]

\textbf{(iii)}
Split into $4$-bit blocks and invert the table:
\[
0101\mapsto F,\quad
0101\mapsto F,\quad
1011\mapsto L,\quad
0110\mapsto G,\quad
0101\mapsto F.
\]
\[
\boxed{\text{FFLGF}}
\]

\textbf{(iv)}
The second block is $1100$, which is $12$ in decimal. Every letter code in the table is an integer from $0$ through $11$, so $1100$ is not the code of any letter.
\[
\boxed{\text{non-decodable}}
\]

\subsection*{(d)}

There are $12^{5}=248832$ strings of length $5$. Since
\[
2^{17}=131072<248832\le 262144=2^{18},
\]
distinguishing them takes $\lceil\log_{2}(12^{5})\rceil=\lceil 5\log_{2} 12\rceil=18$ bits.

\[
\boxed{18}
\]

\subsection*{(e)}

Number the letters as in part (b): $A=0,B=1,\ldots,K=10,L=11$. Encode a string $c_{4}c_{3}c_{2}c_{1}c_{0}$ by the integer
\[
N=c_{4}\cdot 12^{4}+c_{3}\cdot 12^{3}+c_{2}\cdot 12^{2}+c_{1}\cdot 12+c_{0},
\]
write $N$ in binary, and pad on the left with zeros to $18$ bits. The values of $N$ are exactly $0,1,\ldots,12^{5}-1$, and $12^{5}\le 2^{18}$, so every string gets a distinct $18$-bit code.

To decode an $18$-bit string, read it as an integer $N$. If $N\ge 12^{5}$, the string is non-decodable. Otherwise convert $N$ to five base-$12$ digits, using a leading $0$ when there are fewer than five, and replace each digit by its letter.

\textbf{(i)}
$\text{BELCH}$ has digits $1,4,11,2,7$, so
\begin{align*}
N
&= 1\cdot 12^{4}+4\cdot 12^{3}+11\cdot 12^{2}+2\cdot 12+7\\
&= 20736+6912+1584+24+7\\
&= 29263.
\end{align*}
In binary, $29263=111001001001111_{2}$, which is $15$ bits. Three leading zeros make $18$ bits:
\[
\boxed{000111001001001111}
\]

\textbf{(ii)}
$\text{LIKED}$ has digits $11,8,10,4,3$, so
\begin{align*}
N
&= 11\cdot 12^{4}+8\cdot 12^{3}+10\cdot 12^{2}+4\cdot 12+3\\
&= 228096+13824+1440+48+3\\
&= 243411.
\end{align*}
In binary, $243411=111011011011010011_{2}$, already $18$ bits:
\[
\boxed{111011011011010011}
\]

\textbf{(iii)}
\[
100101000111101010_{2}=152042.
\]
Since $152042<248832$, divide by $12$. The remainders, least significant digit first, are
\begin{align*}
152042&=12\cdot 12670+2,\\
12670&=12\cdot 1055+10,\\
1055&=12\cdot 87+11,\\
87&=12\cdot 7+3,\\
7&=12\cdot 0+7.
\end{align*}
Most significant digit first, the digits are $7,3,11,10,2$, which are $H,D,L,K,C$.
\[
\boxed{\text{HDLKC}}
\]

\textbf{(iv)}
\[
001011100010010111_{2}=47255.
\]
Since $47255<248832$, divide by $12$:
\begin{align*}
47255&=12\cdot 3937+11,\\
3937&=12\cdot 328+1,\\
328&=12\cdot 27+4,\\
27&=12\cdot 2+3,\\
2&=12\cdot 0+2.
\end{align*}
Most significant digit first, the digits are $2,3,4,1,11$, which are $C,D,E,B,L$.
\[
\boxed{\text{CDEBL}}
\]

\end{document}
